\subsubsection{多尺度划分与通信感知归核}
\label{subsec:q3_portfolio}
在缓存时序邻域之外，第三问还比较多种划分尺度与归核方式。设 $d(v)$ 为算子在依赖图中的最长前驱链长度，根节点取0。给定目标层带数 $b$，取层带宽度 $h=\lceil(1+\max_v d(v))/b\rceil$，将 $v$ 归入第 $\lfloor d(v)/h\rfloor$ 个层带。每一层带内的弱连通分量作为不可拆的分配单元：忽略边方向后连通的算子保持在一起，层带之间的依赖则只向后延伸。

记分量 $U$ 在流水线 $p$ 上的工作量为 $w_p(U)=\sum_{v\in U:\rho(v)=p}\max\{1,p_v\}$，其中 $p_v$ 为算子给定周期，$\rho(v)$ 为所属流水线。按最大流水线工作量、向上关键性及最小算子标识依次排序分量，再逐一归核。对待分配的 $U$，以
\begin{equation}
 J(U,k)=0.7\max_p\{A_{kp}+w_p(U)\}
       +0.3\max_p\{B_{kp}+w_p(U)\}
       +\lambda\widehat T_{\mathrm{comm}}(U,k),
 \label{eq:q3_band_score}
\end{equation}
\noindent 式中，$A_{kp}$ 为核心 $k$ 在当前层带已分配的流水线工作量，$B_{kp}$ 为该核心各已处理层带的累计工作量；$\widehat T_{\mathrm{comm}}$ 是新增通信的周期估计，$\lambda$ 为其权重。选择 $J$ 最小的核心；并列时依次比较该层带累计通信估计和核心标识。分配后同步更新 $A,B$、张量来源核心及已经计入的通信关系。

通信估计按“张量、来源核心、消费核心”三元组去重。同核生产的数据不重复计入跨核读写，已有同一三元组也不重复收费；外部输入只计读入。张量大小为 $s$ 时，普通读入估计为 $s/\beta_{\mathrm{DDR}}$，跨核数据另加一次 $s/\beta_{\mathrm{DDR}}$ 写出和交接等待。仅当张量可放入缓存、已在较早层带进入模拟队列且仍驻留时，读入估计取 $\tfrac12(s/\beta_{\mathrm{DDR}}+s/\beta_{\mathrm{L2}})$。模拟队列按字节容量先进先出，命中不刷新位置；不把同层带的并发读取预先视为命中。这一折中估计只用于归核排序，实际填充、争用及周期仍由完整事件评价确定。

同层带、同核心的分量合并为一个子图，各核心按层带顺序排列。普通负载归核取目标层带数 $2,4,16,32,64$；通信感知归核取 $(b,\lambda)=(4,0.25),(8,0.25),(16,0.25),(8,0.75),(32,0.25),(64,0.25)$。去除与已有方案重复的结果，并检查合图无环。目标层带数超过依赖深度时可能得到相同划分，故不同参数不必产生不同候选。

另一类候选保持原图的独立弱连通分量完整。归核时同时估计各计算流水线、外部输入读入和输出写出负载，以最大资源负载加其总和的 $0.15$ 倍排序；已经在该核心使用的外部输入不重复计入读入。分量按标识、相反次序或主要共享输入组织核内序列，也可将相邻分量按输入张量并集的字节容量合并。这里的容量约束用于控制输入工作集，不代替实际缓冲分配和FIFO驻留检查。

\subsubsection{归属保持的拆分重排与候选优选}
\label{subsec:q3_selection}
对于已有方案，还可在每个子图内部按弱连通分量拆成单元，保留所有算子的核心归属，再对单元依赖图作拓扑排序。令 $I(U)$ 为单元 $U$ 需要从核外读取的原始张量集合，$\widehat{\mathcal Q}$ 为用于排序的模拟驻留集合，定义
\begin{equation}
 h_U=\sum_{t\in I(U)\cap\widehat{\mathcal Q}}b_t,\qquad
 a_U=\sum_{t\in I(U)\setminus\widehat{\mathcal Q}}b_t,
 \label{eq:q3_unit_bytes}
\end{equation}
\noindent 式中，$b_t$ 为张量字节数，$h_U,a_U$ 分别为估计可复用和需新读入的字节量。就绪单元可按 $(-h_U,a_U)$ 字典序升序选择，或再按向上关键性降序打破并列；小工作集候选则按 $(\mathbf1\{a_U>B_{\mathrm{L2}}\},a_U,-h_U)$ 优先。剩余并列按原父子图位置和最小算子标识确定。每选出一个单元，更新模拟FIFO队列及后继就绪状态；最后将全局拓扑次序投影到各核心。该变换保持归核，但允许改变划分和核内顺序，与算法~\ref{alg:fifo} 的纯顺序移动不同。

这些构造及轨迹窗口候选组成有限候选集。设 $x^0$ 为进入优选阶段的可行方案，设置相对周期改善门槛 $\varepsilon=10^{-3}$。先以候选已有评价周期筛除达不到当前方案改善门槛的结果，再在相同配置下完整确认，并以 $(C_3,D_3)$ 字典序更新。最终方案相对 $x^0$ 也须达到该门槛；这不改变统一迭代中式~\eqref{eq:accept} 的接受规则。
\begin{algorithm}[htbp]
\caption{第三问的多尺度候选优选与同方案配对}\label{alg:q3_portfolio}
\small
\begin{algorithmic}[1]
\Require 可行方案 $x^0$，有限候选集，核数及硬件参数，优选时间上限
\Ensure 正式方案 $x$、完整缓存结果及同方案无L2结果
\State 置 $x\gets x^0$；按构造、通信感知归核及归属保持重排得到候选并去重
\For{每个在剩余时间内可完成评价的候选 $y$}
 \State 若候选已有评价周期大于 $(1-\varepsilon)C_3(x)$，则跳过
 \State 检查覆盖、归核、核内次序及合图无环；不满足则跳过
 \State 在相同硬件配置下完整评价 $y$；未完成或执行不可行则跳过
 \State 若 $(C_3(y),D_3(y))<_{\rm lex}(C_3(x),D_3(x))$，则置 $x\gets y$
\EndFor
\State 若 $[C_3(x^0)-C_3(x)]/C_3(x^0)<\varepsilon$，则恢复 $x\gets x^0$
\State 固定 $x$ 的划分、归核和全部核内顺序，关闭L2作完整配对评价
\State \Return $x$ 及两种缓存条件下的周期、额外COPY和命中字节记录
\end{algorithmic}
\end{algorithm}

有限候选优选与自适应迭代分别受各自时间上限约束。最终性能报告使用选定方案的完整评价，搜索时间曲线则只使用相应搜索阶段在各时刻已经保存的可用解，不将最终优选结果提前计入曲线。
